% ==========================================
% PREÁMBULO DEL DOCUMENTO
% Configuración, paquetes y estilos
% ==========================================

\usepackage[utf8]{inputenc}
\usepackage[english]{babel}
\usepackage{amsmath,amsfonts,amssymb}
\usepackage{booktabs}
\usepackage{graphicx}
\usepackage{multirow}

% --- Configuración de Tipografía ---
% Usamos lmodern en su variante Sans-Serif para un aspecto ultra-moderno, limpio y de vanguardia (Fintech/HFT)
\usepackage{lmodern}
\renewcommand{\familydefault}{\sfdefault}
\usepackage[T1]{fontenc}
\usepackage{microtype} % Mejora el espaciado y justificación

% --- TikZ y PGFPlots ---
\usepackage{tikz}
\usepackage{pgfplots}
\pgfplotsset{compat=1.18}
\usetikzlibrary{shapes.geometric, arrows.meta, positioning, fit, calc, backgrounds, shadows, decorations.pathreplacing}

% Paleta de colores profesional (Inspirada en paletas HFT/Fintech)
\definecolor{techblue}{RGB}{10, 50, 120}
\definecolor{neoncyan}{RGB}{0, 200, 255}
\definecolor{alertred}{RGB}{220, 20, 60}
\definecolor{hotgreen}{RGB}{0, 180, 80}
\definecolor{coldblue}{RGB}{100, 150, 250}

% Configuración global de PGFPlots para gráficos más modernos
\pgfplotsset{
    modern chart/.style={
        ybar,
        bar width=22pt,
        enlarge x limits=0.25,
        width=\linewidth,
        height=6.5cm,
        axis line style={opacity=0}, % Oculta los bordes feos
        major grid style={dashed, gray!30},
        ymajorgrids=true,
        xmajorgrids=false,
        tickwidth=0pt,
        nodes near coords,
        nodes near coords align={vertical},
        nodes near coords style={font=\bfseries\small},
        legend style={
            at={(0.5,-0.2)},
            anchor=north,
            legend columns=-1,
            draw=none, % Sin borde en la leyenda
            fill=white!80
        },
        label style={font=\bfseries},
        tick label style={font=\small}
    }
}

% --- Estilos de Código (Listings) ---
\usepackage{listings}
\usepackage{xcolor}

\definecolor{codegreen}{rgb}{0,0.6,0}
\definecolor{codegray}{rgb}{0.5,0.5,0.5}
\definecolor{codepurple}{rgb}{0.58,0,0.82}
\definecolor{backcolour}{rgb}{0.96,0.96,0.96}

\lstdefinestyle{mystyle}{
    backgroundcolor=\color{backcolour},   
    commentstyle=\color{codegreen}\itshape,
    keywordstyle=\color{techblue}\bfseries,
    numberstyle=\tiny\color{codegray},
    stringstyle=\color{codepurple},
    basicstyle=\ttfamily\footnotesize,
    breakatwhitespace=false,         
    breaklines=true,                 
    captionpos=b,                    
    keepspaces=true,                 
    numbers=left,                    
    numbersep=8pt,                  
    showspaces=false,                
    showstringspaces=false,
    showtabs=false,                  
    tabsize=2,
    frame=leftline,
    framerule=2.5pt,
    rulecolor=\color{techblue!50}
}
\lstset{style=mystyle}

% --- Hyperref ---
\usepackage[colorlinks=true, linkcolor=techblue, citecolor=techblue, urlcolor=techblue]{hyperref}

% --- Balanceo de última página ---
\usepackage{balance}

% --- Espaciado y Ahorro de Espacio para Límite Estricto de 3 Páginas ---
\setlength{\textfloatsep}{6pt plus 2pt minus 2pt}
\setlength{\intextsep}{6pt plus 2pt minus 2pt}
\setlength{\dbltextfloatsep}{6pt plus 2pt minus 2pt}
\setlength{\abovecaptionskip}{3pt plus 1pt minus 1pt}
\setlength{\belowcaptionskip}{2pt plus 1pt minus 1pt}

% Ajuste sutil de listas para evitar saltos huérfanos
\usepackage{enumitem}
\setlist[enumerate]{topsep=2pt, partopsep=0pt, parsep=0pt, itemsep=2pt}
\setlist[itemize]{topsep=2pt, partopsep=0pt, parsep=0pt, itemsep=2pt}
